% Results -- drop-in fragment. Requires \usepackage{booktabs,graphicx}.
% Figures referenced live in experiments/plots/ (copy/point \includegraphics as needed).

\section{Results}
\label{sec:results}

\subsection{Trace-Count Threshold (C)}
\label{sec:results-threshold}
Table~\ref{tab:c-threshold} reports, for the canonical C run, the verified
recovery rate at each trace count with 95\% Wilson score intervals. Recovery
rises from~0\% at $\le 35$k to a reliable~100\% at $\ge 100$k. Fitting a logistic
model in $\log_2(\text{count})$ places the 50\% threshold at
$N_{50} \approx 51{,}346$ traces ($2^{15.65}$); a model-independent linear
interpolation agrees at $\approx 50{,}000$ ($2^{15.61}$). The attack first
reaches high reliability ($\ge 90\%$) at $\approx 75$k traces
(Fig.~\ref{fig:c-success}).

\begin{figure}[t]
\centering
\includegraphics[width=0.85\linewidth]{go_10x_success_rate.png}
\caption{C key-recovery success rate vs.\ trace count (log$_2$ axis) with 95\%
Wilson intervals, the fitted logistic curve, and the 50\% threshold
($N_{50}\approx 51{,}346$).}
\label{fig:c-success}
\end{figure}

\begin{table}[t]
\centering
\caption{C recovery rate vs.\ trace count (canonical run; 10 trials each,
95\% Wilson CI). $N_{50}\approx 51{,}346$.}
\label{tab:c-threshold}
\begin{tabular}{rccc}
\toprule
Traces & Recovered & Rate & 95\% CI \\
\midrule
$\ge 100{,}000$ & 10/10 & 100\% & [72, 100] \\
75{,}000 & 9/10 & 90\% & [60, 98] \\
60{,}000 & 7/10 & 70\% & [40, 89] \\
50{,}000 & 5/10 & 50\% & [24, 76] \\
45{,}000 & 4/10 & 40\% & [17, 69] \\
40{,}000 & 1/10 & 10\% & [2, 40] \\
$\le 35{,}000$ & 0--1/10 & $\le$10\% & [0, 40] \\
\bottomrule
\end{tabular}
\end{table}

Across all counts the mean \texttt{min\_time} stays within 112--121 ticks and
the mean attack score within $\approx 17{,}200$--$17{,}460$ (Fig.~\ref{fig:covar}).
Neither trends with count: the per-trace leak is unchanged, and what improves
with more traces is only the stability of the group-mean estimates. The
threshold is thus a statistical (averaging) effect, not a signal cliff.

\begin{figure}[t]
\centering
\includegraphics[width=\linewidth]{go_10x_covariates.png}
\caption{Supporting evidence: mean \texttt{min\_time} (left) and mean attack
score (right) are flat across trace count, so the per-trace leak is unchanged and
the threshold reflects averaging quality rather than a signal cliff.}
\label{fig:covar}
\end{figure}

\subsection{Cross-Language Recovery}
\label{sec:results-crosslang}
Under identical counts, key, dials, and recovery logic, the three
implementations reach 50\% recovery at very different trace counts
(Table~\ref{tab:crosslang}, Fig.~\ref{fig:crosslang}). Go requires
$N_{50}\approx 265{,}283$ ($2^{18.0}$) with the garbage collector enabled---about
$5.2\times$ the C threshold. Python is deferred (native collection at these
counts is prohibitively slow and expected to recover near~0; see limitations).

\begin{table}[t]
\centering
\caption{Fitted 50\% trace threshold by implementation (identical parameters).}
\label{tab:crosslang}
\begin{tabular}{lccc}
\toprule
Implementation & $N_{50}$ (traces) & $\log_2 N_{50}$ & Ratio to C \\
\midrule
C                 & 51{,}346  & 15.65 & 1.0$\times$ \\
Go (GC on)        & 265{,}283 & 18.02 & 5.2$\times$ \\
Go (GC off)       & 243{,}335 & 17.89 & 4.7$\times$ \\
\bottomrule
\end{tabular}
\end{table}

\begin{figure}[t]
\centering
\includegraphics[width=\linewidth]{cross_language_recovery.png}
\caption{Cross-language recovery. C (blue) reaches 50\% recovery near 50k traces;
Go with the garbage collector on (red) and off (green) require roughly five times
as many, their curves nearly coincident. Points show 95\% Wilson intervals;
dashed lines mark each fitted $N_{50}$.}
\label{fig:crosslang}
\end{figure}

\subsection{Runtime-Noise Ablation (Go GC on vs.\ off)}
\label{sec:results-gc}
Disabling the Go garbage collector and pinning to one core lowers the fitted
threshold only marginally, from $N_{50}\approx 265$k to $\approx 243$k
($\approx 8\%$), well inside the sampling uncertainty. At matched counts the
GC-off run is equal or slightly better everywhere (Table~\ref{tab:gc}), but the
$\approx 5\times$ gap to C is essentially untouched. Quieting the runtime helps
a little; it does not explain Go's disadvantage.

\begin{table}[t]
\centering
\caption{Go recovery at matched counts, garbage collector on vs.\ off
(10 trials each).}
\label{tab:gc}
\begin{tabular}{rcc}
\toprule
Traces & GC on & GC off \\
\midrule
500{,}000 & 8/10 & 9/10 \\
250{,}000 & 4/10 & 6/10 \\
100{,}000 & 0/10 & 0/10 \\
50{,}000  & 0/10 & 0/10 \\
\bottomrule
\end{tabular}
\end{table}

\subsection{Dial Sweep}
\label{sec:results-dial}
At the fixed 100k count, the C dial sweep (Table~\ref{tab:dial}) shows that
increasing \texttt{-repeat} beyond the $r=50$ baseline does \emph{not} improve
recovery and appears to reduce it (70\% at $r=50$ vs.\ 20\% at $r=75,100$),
while the strongest cache disturbance, $e=4096$~KiB, gives the highest recovery
(100\%). The individual points carry wide intervals ($n=10$); the robust reading
is that signal strength (\texttt{-evict-kb}) matters more than additional
averaging (\texttt{-repeat}) at this count (Fig.~\ref{fig:dial}).

The Go dial sweep is uninformative by construction: 100k traces lie far below
Go's $\approx 265$k threshold, so recovery is $\approx 0$ (1/110 with GC on,
1/60 with GC off) regardless of dial setting. This is a consequence of fixing
the dial count at C's knee, which is below Go's operating point---a design
observation, not a null effect (Section~\ref{sec:analysis-dial}).

\begin{table}[t]
\centering
\caption{C dial sweep at 100k traces (10 trials each). Left: \texttt{-repeat}
at $e=2048$. Right: \texttt{-evict-kb} at $r=50$.}
\label{tab:dial}
\begin{tabular}{rc@{\hskip 2em}rc}
\toprule
\texttt{-repeat} & Rate & \texttt{-evict-kb} & Rate \\
\midrule
10  & 60\% & 256  & 80\%  \\
20  & 50\% & 512  & 60\%  \\
30  & 50\% & 1024 & 10\%  \\
50  & 70\% & 2048 & 50\%  \\
75  & 20\% & 4096 & 100\% \\
100 & 20\% &      &       \\
\bottomrule
\end{tabular}
\end{table}

\begin{figure}[t]
\centering
\includegraphics[width=\linewidth]{knee100k_dials.png}
\caption{C dial sweep at 100k traces: recovery rate vs.\ \texttt{-repeat} (left)
and vs.\ \texttt{-evict-kb} (right), with 95\% Wilson intervals. Stronger cache
eviction raises recovery; additional repetition does not.}
\label{fig:dial}
\end{figure}
